% Unpromoted research note, 2026-11-19 after20:37 source freeze.
% No result ID assigned. No change to an included part or ledger is made.
\paragraph{Candidate ordinary TOP sphere formula.}
For $n\ge2$, put
\[
 Y_n=\operatorname{hofib}(B\mathrm O(n)\longrightarrow
                             B\mathrm{TOP}(n)),
\]
where $\mathrm{TOP}(n)=\operatorname{Homeo}(\mathbb R^n;0)$
has its actual compact-open topology. The candidate is a weak
equivalence
\[
 P_{\mathrm{TOP}}(S^n)
   \simeq\operatorname{hofib}_*(Y_n\longrightarrow Y_{n+1}),
                                                               \tag{S.1}
\]
with the arrow the ordinary product-with-a-line stabilization.
All homotopy fibres use the displayed standard basepoints. The
purpose of the proof below is to identify this arrow, not merely
the abstract types of its source and target.

\paragraph{Conventions and opened primary inputs.}
Every topological group below is replaced by its singular simplicial
group. In particular orthogonal groups carry continuous singular
chains, and their inclusion in TOP is literal. Let $\mathcal G_n$
be the simplicial group of germs at $0$ of projection-preserving
homeomorphisms on $\Delta^q\times U$, with one common open
neighbourhood $U$ for the entire simplex. Equality means equality
on a smaller common neighbourhood. These are the finite-parameter
germs of A.132, not block maps.

The current published original of Salvatore--Turchin,
\emph{On the smoothing theory delooping of disc diffeomorphism and
embedding spaces}, Math. Z.313(2026), article30,
DOI10.1007/s00209-026-04025-4, Section3.4, equations(3.5)--(3.7),
explicitly uses this convention. Its full text was opened at
\texttt{https://link.springer.com/article/10.1007/s00209-026-04025-4}.
The primary globalization input it cites was also opened:
Kuiper--Lashof, \emph{Microbundles and bundles I}, Invent. Math.1
(1966), Lemma3, printedp11, with proofpp12--13. That lemma extends
a fibrewise embedding over a parameter simplex, after shrinking its
source, to a fibrewise sphere homeomorphism preserving both marked
points. The relative boundary extension argument immediately after
the lemma is also explicit. Its full original is publicly hosted at
\texttt{https://www.maths.gla.ac.uk/\string~mpowell/Kuiper-Lashof1966\_Article\_MicrobundlesAndBundles.pdf}.
This note does not use the smooth delooping theorem, which excludes
dimension4, from Salvatore--Turchin.

\paragraph{Global and spherical representatives of germs.}
Restriction gives a degreewise surjective simplicial group map
\[
 T_n:=\operatorname{Sing}\operatorname{Homeo}(\mathbb R^n;0)
        \longrightarrow\mathcal G_n.                  \tag{S.2}
\]
Indeed compactness of each simplex gives a common source ball and
a common image chart for its germ representative. Fixed changes of
scale put it into the setting of Kuiper--Lashof Lemma3. The resulting
sphere extension fixes both $0$ and $\infty$; removing $\infty$
gives a global family on $\mathbb R^n$. A degreewise surjective
map of simplicial groups is a Kan fibration.

The kernel of (S.2) consists of families that are the identity on
one common ball about $0$. Inversion in the sphere interchanges
$0$ and $\infty$ and identifies this kernel with compactly
supported families on $\mathbb R^n$ that fix $0$. Such families
are weakly contractible by the Alexander formula
$a_s(x)=s\,a(x/s)$, $s>0$, with $a_0=1$.
Each finite family has common bounded support; this gives uniform
continuity of the map and inverse at $s=0$, and preserves the fixed
point. It follows that (S.2) is a weak equivalence, including its
component assertion.

Fix $p\in S^n$, choose an orthogonally equivariant chart at $p$,
and let $H_n=\operatorname{Sing}\operatorname{Homeo}(S^n;p)$.
The germ map
\[
                        H_n\longrightarrow\mathcal G_n
                                                               \tag{S.3}
\]
is likewise degreewise surjective by the same original lemma.
Its kernel is the group of sphere families that are the identity
on a common neighbourhood of $p$. On $S^n\setminus\{p\}$ it is
the group of compactly supported families. The same Alexander
formula contracts each finite family; there is no additional fixed
point condition in this instance. Hence (S.3) is also a weak
equivalence. These statements use common supports for finite
families; no claim about a union with a different topology is needed.

\paragraph{The first sphere fibre model, before identifying its map.}
Let $A_{n+1}=\operatorname{Sing}\operatorname{Homeo}(D^{n+1};0)$.
The germ restriction
$A_{n+1}\to\mathcal G_{n+1}$ is a Kan fibration. The proof in
A.132 is dimension-independent: the Edwards--Kirby local extension
operator is fixed on a neighbourhood of the identity; a common
source/image ball and a single finite mesh for each compact
parameter family extend successive increments, keeping prescribed
horns exactly fixed. No globalization of a disk vertex is needed
to prove this fibration assertion.

Its identity fibre is the filtered simplicial union of
$\operatorname{Sing}\operatorname{Homeo}(D^{n+1}\ \mathrm{rel}\
D_\epsilon^{n+1})$. Each stage is the ordinary free-outer-boundary
concordance space of $S^n$. The radial conjugation and compression
argument in A.132 applies in the same dimension-independent way:
any finite parameter family and any relative nullhomotopy admit a
common $\epsilon$, and inclusions of stages are weak equivalences.
Consequently the fibre is weakly equivalent to
$P_{\mathrm{TOP}}(S^n)$, with the inner face fixed and the outer
face free.

Boundary restriction $A_{n+1}\to
G_n:=\operatorname{Sing}\operatorname{Homeo}(S^n)$ is a weak
equivalence. It has the radial cone section, and its kernel consists
of boundary-fixed disk homeomorphisms also fixing $0$, contracted
by the relative Alexander formula. Extend every disk map outside
the disk by its coned boundary restriction. This gives a group map
$A_{n+1}\to T_{n+1}$ that commutes with germ restriction.
For the cone section its composite is the actual radial cone map
\[
 c:G_n\longrightarrow T_{n+1},\qquad
 c(h)(r\omega)=r h(\omega).                             \tag{S.4}
\]
Together with (S.2), these weak equivalences give
\[
 P_{\mathrm{TOP}}(S^n)\simeq\operatorname{hofib}_1(c).
                                                               \tag{S.5}
\]
The full germ equivalence (S.2) permits all components in this
formula; the proof is not restricted to an unproved identity-component
replacement of $G_n$.

\paragraph{Orthogonal quotients and the point stabilizer.}
The map (S.4) restricts to the standard linear inclusion on
$\mathrm O(n+1)$. Quotienting both its source and target by this
same subgroup therefore preserves its based homotopy fibre:
\[
 \operatorname{hofib}(c)\simeq
 \operatorname{hofib}\bigl(
 G_n/\mathrm O(n+1)\longrightarrow
 T_{n+1}/\mathrm O(n+1)\bigr).                          \tag{S.6}
\]
This follows from the map of principal orthogonal fibrations,
whose map on fibres is the identity. Quotients here can equivalently
be taken in simplicial groups or by homotopy orbits: the action is
free and the orthogonal principal bundles are locally trivial.

The inclusion of the point stabilizer induces an equivalence
\[
 H_n/\mathrm O(n)\longrightarrow G_n/\mathrm O(n+1).
                                                               \tag{S.7}
\]
For actual topological quotients it is a homeomorphism. Every right
coset has a representative fixing $p$, because $\mathrm O(n+1)$
acts transitively on $S^n$, and two such representatives differ
by exactly $\mathrm O(n)$. Local continuous choices of a rotation
taking $p$ to a nearby point provide the local inverses. Passing
to singular complexes gives (S.7); a principal orthogonal bundle
over a simplex has a section, so the simplicial and topological
quotient conventions agree up to the indicated equivalence.

Choose the chart at $p$ to be $\mathrm O(n)$-equivariant. Then
(S.3) is exactly the identity on that orthogonal subgroup and
gives a weak equivalence
\[
 H_n/\mathrm O(n)\xrightarrow{\simeq}
                \mathcal G_n/\mathrm O(n)\simeq Y_n.   \tag{S.8}
\]
Similarly (S.2) identifies the target quotient in (S.6) with
$Y_{n+1}$. It remains essential to identify the arrow between them.

\paragraph{Radial extension gives the actual stabilization map.}
Take $p=(0,\ldots,0,1)$ and $h\in\operatorname{Homeo}(S^n;p)$.
Its cone $c(h)$ fixes every point of the ray through $p$. Therefore
\[
 F_s(h)=\tau_{-sp}\,c(h)\,\tau_{sp},\qquad0\le s\le1,
                                                               \tag{S.9}
\]
is a continuous homotopy of group homomorphisms from $H_n$ to
$T_{n+1}$; every map fixes $0$. For $A\in\mathrm O(n)$,
embedded as $\operatorname{diag}(A,1)$, one has
$F_s(A)=\operatorname{diag}(A,1)$ for every $s$. Thus this homotopy
is exact on the orthogonal subgroup and descends to the quotients.
It moves the germ of $c(h)$ at $0$ to its germ at $p$.

Use the local chart
\[
 \theta(u)=(u,\sqrt{1-\|u\|^2}),\qquad
 \psi(u,t)=(1+t)\theta(u)-p.                            \tag{S.10}
\]
It is defined near $(u,t)=0$, has derivative the identity there,
and commutes with $\operatorname{diag}(A,1)$ for every
$A\in\mathrm O(n)$. Exact substitution gives
\[
 \psi^{-1}F_1(h)\psi(u,t)
                    =(\theta^{-1}h\theta(u),t).       \tag{S.11}
\]
The right-hand side is precisely product with a line on the
point germ of $h$.

Finally the coordinate change in (S.11) introduces no additional
map. The germs $\psi_r(z)=r^{-1}\psi(rz)$ for $r>0$, completed
at $r=0$ by $D\psi(0)=1$, form a continuous family of local
diffeomorphisms and their inverses on a common neighbourhood.
They commute with $\mathrm O(n)$ throughout. Conjugating by this
fixed family gives a homotopy, relative to that subgroup, from
the coordinate-conjugated germ in (S.11) to the original germ.
For any finite parameter family of $h$ a common smaller chart
works, by compactness. The same formulas preserve parameter faces.
Thus under (S.7)--(S.8) the map in (S.6) is the ordinary stabilization
$Y_n\to Y_{n+1}$ up to the specified based homotopy. This proves
the candidate formula (S.1).

\paragraph{Safe low-degree consequence for the four-sphere.}
The corrected Freedman--Quinn statement, audited in A.132 using
Gabai--Gay--Hartman--Krushkal--Powell2026, says that
$Y_4\to\mathrm{TOP}/\mathrm O$ is5-connected. Hence it is an
isomorphism on $\pi_i$ for $i\le4$. The finite-rank statement
Kirby--Siebenmann V, Theorem5.5(I), p251, gives
$\mathrm{TOP}_5/\mathrm{PL}_5\simeq K(\mathbb Z/2,3)$.
Together with the already audited
$\pi_i(\mathrm{PL}_5/\mathrm O_5)=0$ for $i<5$, this shows that
$Y_5\to\mathrm{TOP}/\mathrm O$ is an isomorphism for $i\le4$.
Naturality implies the same for $Y_4\to Y_5$.
The fibre sequence in (S.1) therefore gives
\[
                \pi_iP_{\mathrm{TOP}}(S^4)=0
                           \quad(0\le i\le3).         \tag{S.12}
\]
For PL the audited fibre formula A.128 and the corresponding
low-degree vanishing of $\mathrm{PL}_4/\mathrm O_4$ and
$\mathrm{PL}_5/\mathrm O_5$ give the same conclusion.
No vanishing of $\pi_4$ is inferred: that would require control
of $\pi_5Y_4\to\pi_5Y_5$, which the cited inputs do not supply.
No compatibility of an arbitrary model of PL/orthogonal groups
with continuous orthogonal singular chains is silently assumed.

\paragraph{Scope.}
This is a sphere model and a special low-degree consequence. It is
not a computation of the general remaining low-dimensional
concordance groups, and by itself does not reduce the main general
comparison gate. The finite-rank TOP groups have not been replaced
by their stable group.
